\head{第十九版：第一篇8.14至第11节续校}
\noindent\textbf{本版性质。}以第十八版已公开提交为基线，对第五卷预备手稿第一篇8.14至第11节进行法英逐段回校；这部分旧译原已存在。本轮增加的是底本辨读、遗漏文字的恢复、实质订正及独立核查，不把它们计为从无到有的全卷新译。所有Dot原件及四十批已整合模块仍在工程中，较早各卷正文和本篇第12节以后译文继续保留。

\section*{真实接合边界}
旧输入模块\texttt{v18\_first\_8\_14\_onward.tex}的范围是8.14至第11节。本轮将其替换为五个按原顺序接续的模块：8.14--8.18；8.18之后至第8节末；第9节；第10节；第11节。旧模块原件仍在工程中，未重复输入。下一实际回校点是法稿物理64页第12节，对读英文物理43页中段；第12节旧译在\texttt{v\_dot\_middle.tex}中。

\section*{本轮主要发现}
\begin{enumerate}
\item 8.16的原弱非分歧条件存在余维二收缩反例；正文的加强版保留深度、几何正规、几何整和固有等全部原条件，并明确加强非分歧性。
\item 8.18不能把“不可约但非约化”的特殊截面当作几何整纤维开集中的点。正文采用明确的充分条件版本，原充要说法及证明完整保存于对应报告。
\item 一般线性系必须控制固定分支。投影像二维并不自动保证整个泛截面不可约；包含中心的线性空间维数也须从原法稿的“多一维”改为“多两维”。
\item 9.3的析取词可从法稿恢复，但原命题另缺固定分支条件。正文采用有证明的附加条件，并给爆破反例。9.4、9.5仍为未编成的提示和残文。
\item 第10节恢复闭包、纤维积、零点概形、不等号和剩余域等读法；正则性的一般约定及存在性依赖分别保留。本轮给出局部诺特范围内的爆破补证。
\item 第11节法稿中原有而英译漏掉的完整条件得到恢复；无条件拟紧的断言仍有反例。旧报告若把英译漏项归为法稿缺条件，本轮明确更正归责。
\end{enumerate}

\section*{证据、历史与未决范围}
来源记录标明法稿物理页、手写页和英译物理页；页面图像与散列保存在本版\texttt{recollation/}。字迹恢复只依据实际可读原页；数学上的修正或加强单独标为译校结论。难辨的短词、密集删划、9.5残文及8.8未取得的完整起首不会因新编号而变成已经解决的问题。

本报告后部完整保留第十八版及更早报告，其旧“未见原稿”判断属于历史。新判断仅覆盖本轮直接对读的范围；没有对全部历史材料再作一次数学审定。正式第IV章主体仍未连续译入，第0章与预备手稿章号不能替代它。

\section*{技术验证的记录方法}
以真实\texttt{ctexbook}类排版，不另设正文版心或字体。实际运行记录分别保存在\path{checks/current_structure_verification.json}、\path{checks/current_build_log_summary.json}、\path{checks/current_pdf_parity.json}和\path{checks/visual_verification_summary.json}。它们分别报告输入结构和继承文件、四份最终构建日志、模块与单文件逐页比较、实际查看过的代表页；各项必须有执行结果，不能只凭本段文字判为通过。页数、字节与散列以成品记录\texttt{ARTIFACTS.json}为准。

下面各组报告给出具体原页、订正、反例、补证及未决项。所有补证均为译者文字，不冒作作者被遗漏的证明。
